% Algebraic curves compared: equation (left), real locus (middle), and the
% Riemann surface of the projective closure (right), for an ellipse (genus 0),
% two elliptic curves (genus 1), the nodal and cuspidal cubics, and the Fermat
% quartic (genus 3). Points at infinity are labelled ∞; singular points are red.
% Real loci are drawn from rational or root-anchored parametrizations so each
% branch is one smooth plot; the y-axis is compressed by the factor shown.
\documentclass[tikz,border=2pt]{standalone}
\usepackage{dzg-tikz}

\begin{document}
\begin{tikzpicture}[every label/.append style={inner sep=1pt}]

% Row 1: ellipse (degree 2, genus 0), two points at infinity.
\node[font=\small] at (0,0) {$x^2/a^2 + y^2/b^2 = 1$};
\begin{scope}[shift={(4.8,0)}, scale=0.75]
  \draw[curve component] (0,0) ellipse (1.8 and 1.2);
  \node[lattice point] at (0,0) {};
\end{scope}
\begin{scope}[shift={(10,0)}, scale=0.75]
  \draw[curve component] (0,0) circle (1.5);
  \draw[curve component] (-1.5,0) arc[start angle=180, end angle=360, x radius=1.5, y radius=0.55];
  \draw[curve component, dashed] (1.5,0) arc[start angle=0, end angle=180, x radius=1.5, y radius=0.55];
  \node[curve node, label={above left:$\infty_1$}] at (-1.15,0.2) {};
  \node[curve node, label={above right:$\infty_2$}] at (1.15,0.2) {};
\end{scope}

% Row 2: y^2 = x^3 + x^2 + 1, one real component; its real root is
% r = -1.46557, and x = r + t^2 parametrizes the component.
\begin{scope}[shift={(0,-3.6)}]
  \node[font=\small] at (0,0) {$y^2 = x^3 + x^2 + 1$};
  \begin{scope}[shift={(4.8,0)}, scale=0.75]
    \clip (-2.4,-1.9) rectangle (2.8,1.9);
    \draw[coordinate axis] (-2.3,0) -- (2.5,0);
    \draw[coordinate axis] (0,-1.8) -- (0,1.8);
    \draw[curve component, samples=150, variable=\t, domain=-1.9:1.9, smooth]
      plot ({-1.46557+\t*\t},
            {0.6*sign(\t)*sqrt(max(0, (-1.46557+\t*\t)^3 + (-1.46557+\t*\t)^2 + 1))});
    \node[lattice point] at (0,0) {};
  \end{scope}
  \begin{scope}[shift={(10,0)}, scale=0.75]
    \draw[curve component] (0,0) ellipse (2 and 1.3);
    \draw[curve component] (-0.9,0.1) to[curve through={(0,-0.25)}] (0.9,0.1);
    \draw[curve component] (-0.65,-0.05) to[curve through={(0,0.2)}] (0.65,-0.05);
    \node[curve node, label={right:$\infty$}] at (2.1,0) {};
  \end{scope}
\end{scope}

% Row 3: y^2 = x^3 - x, two real components: the oval over [-1,0] and the
% branch x = 1 + t^2.
\begin{scope}[shift={(0,-7.2)}]
  \node[font=\small] at (0,0) {$y^2 = x^3 - x$};
  \begin{scope}[shift={(4.8,0)}, scale=0.75]
    \clip (-2.4,-1.9) rectangle (2.8,1.9);
    \draw[coordinate axis] (-2.3,0) -- (2.7,0);
    \draw[coordinate axis] (0,-1.8) -- (0,1.8);
    \draw[curve component, samples=150, variable=\t, domain=0:360, smooth]
      plot ({-0.5-0.5*cos(\t)},
            {0.7*sign(sin(\t))*sqrt(max(0, (-0.5-0.5*cos(\t))^3 + 0.5 + 0.5*cos(\t)))});
    \draw[curve component, samples=150, variable=\t, domain=-1.225:1.225, smooth]
      plot ({1+\t*\t}, {0.7*sign(\t)*sqrt(max(0, (1+\t*\t)^3 - 1 - \t*\t))});
    \node[lattice point] at (0,0) {};
  \end{scope}
  \begin{scope}[shift={(10,0)}, scale=0.75]
    \draw[curve component] (0,0) ellipse (2 and 1.3);
    \draw[curve component] (-0.9,0.1) to[curve through={(0,-0.25)}] (0.9,0.1);
    \draw[curve component] (-0.65,-0.05) to[curve through={(0,0.2)}] (0.65,-0.05);
    \node[curve node, label={right:$\infty$}] at (2.1,0) {};
  \end{scope}
\end{scope}

% Row 4: nodal cubic y^2 = x^3 + x^2 = x^2(x+1), parametrized by
% (t^2 - 1, t(t^2 - 1)); its Riemann surface is a torus pinched at the node.
\begin{scope}[shift={(0,-10.8)}]
  \node[font=\small] at (0,0) {$y^2 = x^3 + x^2$};
  \begin{scope}[shift={(4.8,0)}, scale=0.75]
    \clip (-2.4,-1.9) rectangle (2.8,1.9);
    \draw[coordinate axis] (-2.3,0) -- (2.7,0);
    \draw[coordinate axis] (0,-1.8) -- (0,1.8);
    \draw[curve component, samples=150, variable=\t, domain=-1.817:1.817, smooth]
      plot ({\t*\t-1}, {0.7*\t*(\t*\t-1)});
    \node[curve node, fill=dzg singular] at (0,0) {};
  \end{scope}
  \begin{scope}[shift={(10,0)}, scale=0.75]
    \coordinate (pinch) at (0,-1.3);
    \draw[curve component] (pinch)
      to[curve through={(1.5,-0.9) .. (2,0) .. (1.2,1.1) .. (0,1.3) .. (-1.2,1.1) .. (-2,0) .. (-1.5,-0.9)}] (pinch);
    \draw[curve component] (pinch)
      to[curve through={(0.55,-0.35) .. (0,0.35) .. (-0.55,-0.35)}] (pinch);
    \node[curve node, fill=dzg singular] at (pinch) {};
    \node[curve node, label={right:$\infty$}] at (2.1,0) {};
  \end{scope}
\end{scope}

% Row 5: cuspidal cubic y^2 = x^3, parametrized by (t^2, t^3); its Riemann
% surface is a sphere with a cusp.
\begin{scope}[shift={(0,-14.4)}]
  \node[font=\small] at (0,0) {$y^2 = x^3$};
  \begin{scope}[shift={(4.8,0)}, scale=0.75]
    \clip (-2.4,-1.9) rectangle (2.8,1.9);
    \draw[coordinate axis] (-1.5,0) -- (2.7,0);
    \draw[coordinate axis] (0,-1.8) -- (0,1.8);
    \draw[curve component, samples=150, variable=\t, domain=-1.517:1.517, smooth]
      plot ({\t*\t}, {0.7*\t*\t*\t});
    \node[curve node, fill=dzg singular] at (0,0) {};
  \end{scope}
  \begin{scope}[shift={(10,0)}, scale=0.75]
    \coordinate (cusp) at (-1.8,0);
    \draw[curve component] (cusp)
      to[curve through={(-0.6,1.1) .. (0.8,1.2) .. (1.5,0) .. (0.8,-1.2) .. (-0.6,-1.1)}] (cusp);
    \draw[curve component] (cusp) to[curve through={(0,-0.3)}] (1.5,0);
    \draw[curve component, dashed] (cusp) to[curve through={(0,0.3)}] (1.5,0);
    \node[curve node, fill=dzg singular] at (cusp) {};
    \node[curve node, label={right:$\infty$}] at (1.65,0) {};
  \end{scope}
\end{scope}

% Row 6: Fermat quartic x^4 + y^4 = 1 (genus 3), four points at infinity.
\begin{scope}[shift={(0,-18)}]
  \node[font=\small] at (0,0) {$x^4 + y^4 = 1$};
  \begin{scope}[shift={(4.8,0)}, scale=0.75]
    \clip (-2.4,-1.9) rectangle (2.8,1.9);
    \draw[coordinate axis] (-2.3,0) -- (2.3,0);
    \draw[coordinate axis] (0,-1.8) -- (0,1.8);
    \draw[curve component, samples=200, variable=\t, domain=0:360, smooth]
      plot ({1.8*sign(cos(\t))*sqrt(abs(cos(\t)))},
            {1.2*sign(sin(\t))*sqrt(abs(sin(\t)))});
    \node[lattice point] at (0,0) {};
  \end{scope}
  \begin{scope}[shift={(10,0)}, scale=0.75]
    \draw[curve component, use Hobby shortcut, closed=true]
      (-3.5,0) .. (-3.2,1.0) .. (-2.5,1.2) .. (-1.25,0.6) .. (0,1.2) .. (1.25,0.6)
      .. (2.5,1.2) .. (3.2,1.0) .. (3.5,0) .. (3.2,-1.0) .. (2.5,-1.2) .. (1.25,-0.6)
      .. (0,-1.2) .. (-1.25,-0.6) .. (-2.5,-1.2) .. (-3.2,-1.0);
    \foreach \c in {-2.5, 0, 2.5} {
      \draw[curve component] ({\c-0.45},0.05) to[curve through={(\c,-0.2)}] ({\c+0.45},0.05);
      \draw[curve component] ({\c-0.3},-0.07) to[curve through={(\c,0.12)}] ({\c+0.3},-0.07);
    }
    \node[curve node, label={above:$\infty_1$}] at (-1.25,0.6) {};
    \node[curve node, label={above:$\infty_2$}] at (1.25,0.6) {};
    \node[curve node, label={below:$\infty_3$}] at (-1.25,-0.6) {};
    \node[curve node, label={below:$\infty_4$}] at (1.25,-0.6) {};
  \end{scope}
\end{scope}

\end{tikzpicture}
\end{document}
